\section{1st modular law for area variations}

Let us define the two normalized areas
\begin{equation}
\label{rm:eq:modular-normalized-area-channels}
\mathcal A_m=\gammaH a_0N_m,
\qquad m\in\{90,120\}.
\end{equation}
Then
\begin{equation}
\label{rm:eq:modular-beta-definition}
\HHM^{(A)}-cI
=\beta_{90}\mathcal A_{90}+\beta_{120}\mathcal A_{120},
\qquad
\beta_m=\frac{m\Dmod}{\gammaH a_0}.
\end{equation}
The terminal values are
\begin{align}
\beta_{90}&=5.9726485401070929914651798889\ldots,
\label{rm:eq:modular-beta90}\\
\beta_{120}&=7.9635313868094573219535731852\ldots,
\qquad
\frac{\beta_{120}}{\beta_{90}}=\frac43.
\label{rm:eq:modular-beta120}
\end{align}

\begin{theorem}[First modular law in areal variables]
\label{rm:thm:modular-first-law}
For a differentiable family of normalized states
\(\rho(\varepsilon)\) passing through \(\rho_A\), while keeping
\(\Dmod,\gammaH,a_0\) fixed,
\begin{equation}
\label{rm:eq:modular-first-law}
\boxed{
\delta S
=\beta_{90}\,\delta\langle\mathcal A_{90}\rangle
+\beta_{120}\,\delta\langle\mathcal A_{120}\rangle.}
\end{equation}
\end{theorem}

\begin{proof}
The first entanglement law in the reference state is \(\delta S=\delta\langle-\log\rho_A\rangle\). The
scalar term \(cI\) does not contribute because
\(\delta\Tr\rho=0\). Substituting \cref{rm:eq:modular-beta-definition} proves the identity.
\end{proof}

The formula is local in state space. For finite changes,
relative entropy appears; omitting it would turn a tangent identity into a
false claim of global linearity.

\Needspace{22\baselineskip}
\begin{theorem}[Finite Identity with a Defect Term]
\label{rm:thm:modular-finite-first-law}
For every state \(\sigma\) whose support is contained in that of
\(\rho_A\),
\begin{equation}
\label{rm:eq:modular-finite-first-law}
\boxed{
S(\sigma)-S(\rho_A)
=\beta_{90}\Delta_\sigma\langle\mathcal A_{90}\rangle
+\beta_{120}\Delta_\sigma\langle\mathcal A_{120}\rangle
-D(\sigma\Vert\rho_A),}
\end{equation}
where
\(D(\sigma\Vert\rho_A)=\Tr\sigma(\log\sigma-\log\rho_A)\geq0\).
In particular, the linear variation of
\cref{rm:thm:modular-first-law} is the tangent term of this identity.
\end{theorem}

\begin{proof}
By definition and by \(\HHM^{(A)}=-\log\rho_A\),
\[
D(\sigma\Vert\rho_A)
=-S(\sigma)+\Tr(\sigma\HHM^{(A)}).
\]
For \(\sigma=\rho_A\), the left-hand side is zero and
\(S(\rho_A)=\Tr(\rho_A\HHM^{(A)})\). Subtracting the two identities gives
\[
S(\sigma)-S(\rho_A)
=\Delta_\sigma\langle\HHM^{(A)}\rangle
-D(\sigma\Vert\rho_A).
\]
The scalar \(cI\) cancels between normalized states; substituting
\cref{rm:eq:modular-beta-definition} yields
\cref{rm:eq:modular-finite-first-law}. The positivity of relative entropy
is Klein's inequality.
\end{proof}

The correction \(D(\sigma\Vert\rho_A)\) measures the informational distance
between the actual state and the reference modular state. Thus, the first
law loses no content upon leaving the infinitesimal regime: it becomes
an exact identity with a controlled positive defect.
